\summarypart{FIVE}{The Age of the Reformation}{1517--1648}
\begin{partoverview}
The reform that the late Middle Ages had demanded arrived suddenly in the sixteenth century. It ended
the pope's hold on the Western church and produced Protestantism in four chief branches: Lutheran,
Reformed, Anabaptist and Anglican. Rome answered with a renewal of its own, in which the new Society of
Jesus took the lead. Then came a century of religious wars. Only after they had exhausted Europe did the
two sides accept that the old unity would not return, and a handful of thinkers began to describe the
church in a new way, as one body existing in several denominations.
\end{partoverview}
\partchapters{%
  \pch{24}{A Wild Boar in the Vineyard}{Martin Luther and Protestantism}
  \pch{25}{Radical Discipleship}{The Anabaptists}
  \pch{26}{Thrust into the Game}{John Calvin}
  \pch{27}{The Curse upon the Crown}{The Church of England}
  \pch{28}{``Another Man'' at Manresa}{The Catholic Reformation}
  \pch{29}{Opening the Rock}{America and Asia}
  \pch{30}{The Rule of the Saints}{Puritanism}
  \pch{31}{Unwilling to Die for an Old Idea}{Denominations}}
